% ===================================================================
%  اللوحة رقم ٣ — الطفولة والشباب.
%
%  Traduction, faite en 2026, en arabe levantin d'aujourd'hui, sur
%  l'ido de texto/io. Regles de la colonne : voir 10-lawha-01.tex.
%  Deux scenes, deux numerotations.
%
%  LE LEVANT A GARDE UN CALENDRIER PLUS VIEUX QUE L'ARABE, ET CE
%  TABLEAU EST L'ENDROIT OU ON LE VOIT. L'alinea 1 de la deuxieme
%  scene nomme trois mois d'ete. L'Egypte et le Golfe disent يوليو،
%  أغسطس، سبتمبر, transcrits du latin par l'italien. Le Levant — la
%  Syrie, le Liban, la Palestine, la Jordanie, et l'Irak avec eux —
%  dit « حزيران »، « تموز »، « آب », et ainsi de suite pour les douze :
%  كانون، شباط، آذار، نيسان، أيار، أيلول، تشرين. CE SONT DES NOMS
%  AKKADIENS, arrives par le syriaque, et تموز est Tammuz, le dieu
%  qu'on pleurait a Byblos quand l'Adonis coulait rouge au printemps.
%  La colonne haoussa avait emprunte ses douze mois a l'anglais avec
%  la chose ; celle-ci ecrit des mots qui ont deux mille ans de plus
%  que la langue dans laquelle ils sont ecrits.
%
%  LE TRIPLET DE korekti.json ETAIT POSE AVANT QUE CE TABLEAU SOIT
%  ECRIT. Le fac-simile ido compose « \VUgras{junio} (julio od
%  \VUgras{agosto}) » et oublie le gras sur le mois du milieu ;
%  gravuri/korekti.json le retablit, colonne par colonne, et la ligne
%  « apc » y a ete deposee au moment du cablage. Elle attendait ce
%  fichier-ci : « (تموز أو » devient « (<b>تموز</b> أو ».
%
%  ET LA SEMAINE DIT LA MEME CHOSE QUE LE CORPS DU TABLEAU 2. Le
%  lundi est « التنين » et non الاثنين, le mardi « التلاتا » et non
%  الثلاثاء : le ث est tombe dans le ت, exactement comme le ظ etait
%  tombe dans le ض. Le taureau (45) est « التور » et non الثور, les
%  vetements (4) sont « التياب » et non الثياب. UNE SEULE CONSONNE, ET
%  ELLE TRAVERSE LE TABLEAU DU LUNDI JUSQU'AUX ARENES.
%
%  LES DEUX ARBRES DU VILLAGE SONT D'ICI, ET C'EST LA PREMIERE FOIS
%  QUE LE RELEVE PEUT L'ECRIRE SANS DETOUR. Le peuplier (3) est
%  « الحور », le platane (17) est « الدلب » : les deux poussent dans
%  la Bekaa et le long du Barada, les deux ont leurs noms depuis
%  toujours. La colonne gujaratie avait du transcrire le sapin faute
%  d'exemplaire ; ici il n'y a rien a transcrire.
%
%  Trois mots turcs de plus, apres la صوبيا et le بوري du tableau 1 :
%  la vieille femme (71) est « ختيارة », de ihtiyar, le parapet (27)
%  est « الدرابزين », et les dragees (55) sont « ملبّس » — celles-la
%  arabes, mais c'est bien ce qu'on jette a un bapteme levantin, et
%  le livret grave exactement cette scene. Et l'argent (52) est
%  « مصاري », qui est l'egyptien مصر, l'Egypte, passe au pluriel : la
%  piastre d'Egypte a fini par nommer l'argent tout court.
% ===================================================================

%%K t03-apar-1 apar
\VUsaut{3.0mm}
\VUtitre{15.0pt}{60}{اللوحة رقم ٣}
\VUsaut{1.6mm}
\VUtitre{11.4pt}{0}{الطفولة والشباب.}
\VUsaut{3.4mm}

%%K t03-apar-2 apar
(اللوحتين ٣ و٤ مركّبتين هيك لتقدّموا بأربع \VUgras{مشاهد عيلية}
المفاهيم الأساسية عن \VUgras{الطقس} و\VUgras{العمر} و\VUgras{العيلة}
و\VUgras{اللبس} و\VUgras{الحال}.)

\VUsaut{4.0mm}
%%K t03-tit-1 sub
\VUcentre{12.0pt}{0}{\textit{المشهد الأول.}}
\VUsaut{3.0mm}

%%K t03-tit-2 sub
\VUpk{10.2pt}{40}{عمادة بضيعة.}
\VUsaut{2.4mm}

%%K t03-tit-3 sub
\VUpk{10.2pt}{40}{الربيع.}
\VUsaut{3.0mm}

%%K t03-c1-01-1 p
1. --- نحنا بـ\VUgras{الربيع}، بشهر \VUgras{آذار} أو \VUgras{نيسان} أو
\VUgras{أيار}، بأيام \VUgras{الأسبوع}، \VUgras{التنين} أو
\VUgras{التلاتا} أو \VUgras{الأربعا}. \VUgras{الساعة عشرة} الصبح.

%%K t03-c1-02-1 p
2. --- \VUgras{الطقس} حلو و\VUgras{الهوا} نقي. الشمس عم تشرق. وكم
\VUgras{غيمة} \textsuperscript{(2)} صغيرة عم تطلع بـ\VUgras{السما}
\textsuperscript{(1)} الزرقا.

%%K t03-c1-03-1 p
3. --- \VUgras{الشجر} بالكاد عليه \VUgras{ورق}. \VUgras{الحور}
\textsuperscript{(3)} الرفيع و\VUgras{الدلب} \textsuperscript{(17)} العالي
لسّا ما عليهن غير براعم.

%%K t03-c1-04-1 p
4. --- \VUgras{السنونو} \textsuperscript{(4)} رجعوا وعم يطيروا وعم
يزقزقوا بفرح حوالين \VUgras{السهم} \textsuperscript{(7)} تبع \VUgras{برج
الأجراس} \textsuperscript{(5)} اللي بيتوّج \VUgras{كنيسة}
\textsuperscript{(6)} الضيعة.

%%K t03-tit-4 sub
\VUsaut{3.4mm}
\VUpk{10.2pt}{40}{الكنيسة.}
\VUsaut{3.0mm}

%%K t03-c1-05-1 p
5. --- هالكنيسة كتير بسيطة، مع \VUgras{بابها الكبير}
\textsuperscript{(13)} و\VUgras{ساعتها} \textsuperscript{(8)} الكبيرة.
\VUgras{العقرب الزغير} \textsuperscript{(9)} على رقم X، عم يدلّ على
\VUgras{الساعات}. و\VUgras{الكبير} \textsuperscript{(10)} على XII (نص
النهار) ؛ عم يدلّ على \VUgras{الدقايق}.

%%K t03-c1-06-1 p
6. --- ببرج الأجراس، \VUgras{الأجراس} \textsuperscript{(11)} عم تدقّ
بفرح. وعلى قمّة السقف واقف \VUgras{صليب} \textsuperscript{(12)} زغير من
حجر.

%%K t03-c1-07-1 p
7. --- الكنيسة ما إلها بس \VUgras{الباب الكبير} \textsuperscript{(13)}،
إلها كمان باب زغير عالجنب. من هالباب الزغير بيطلع \VUgras{الخوري}
\textsuperscript{(15)} لابس \VUgras{جبّته}، بيطلع من \VUgras{غرفة
الثياب} \textsuperscript{(14)} وبيروح على \VUgras{بيته}
\textsuperscript{(19)}.

%%K t03-c1-08-1 p
8. --- وجنبه، هيّي \VUgras{اللبّانة} \textsuperscript{(24)} مع
\VUgras{حمارها} \textsuperscript{(23)}. راجعة من المدينة، اللي حملت
عليها حليب منيح للولاد.

%%K t03-tit-5 sub
\VUsaut{4.0mm}
\VUpk{10.2pt}{40}{بستان الفواكه.}
\VUsaut{3.4mm}

%%K t03-c1-09-1 p
9. --- سيّد أليبورون (الحمار) كتير مبسوط من هالركضة الصبحية. عم يشمّ
بلذّة الهوا المعطّر بريحة \VUgras{الزهر} \textsuperscript{(20)} اللي
مغطّي \VUgras{شجر الفواكه} \textsuperscript{(18)} تبع \VUgras{البستان}
المجاور. كله وردي وأبيض هالـ\VUgras{شجر الدرّاق} \textsuperscript{(18)}
و\VUgras{شجر المشمش} \textsuperscript{(19)}، وعم يوعدوا بموسم منيح.

%%K t03-c1-10-1 p
10. --- وبنفس الوقت، \VUgras{النحلات} \textsuperscript{(22)} الزغار
تبع \VUgras{الخليّة} \textsuperscript{(21)} عم يروحوا لهونيك يسرقوا
\VUgras{عسلهن} الحلو وهنّي عم يطنّوا.

%%K t03-c1-11-1 p
11. --- بس شو عم يصير ؟ هيّن ولاد الضيعة عم يركضوا وعم يهرّبوا
\VUgras{الوز} \textsuperscript{(25)} المذعور. هاد طلوع الموكب من
الكنيسة بعد \VUgras{عمادة} بنت الكاتب بالعدل.

%%K t03-c1-12-1 p
12. --- ويعقوب الزغير، بـ\VUgras{مهده} \textsuperscript{(26)}، فاق من
دقّ الأجراس المفرح، وأخته مجدلينا، اللي كمان بدّها تشوف موكب العمادة،
عم تترجّى إمّها تيجي تمشّطها وتلبّسها عند عتبة البيت.

%%K t03-c1-13-1 p
13. --- الإم وافقت. حطّت \VUgras{منديلها} \textsuperscript{(28)}
و\VUgras{صدريتها} \textsuperscript{(29)} و\VUgras{مريولها}
\textsuperscript{(30)}، وإجت قعدت على كرسي زغير قدّام الباب. ومجدلينا
ما عليها غير \VUgras{تنّورتها الداخلية} \textsuperscript{(31)}
و\VUgras{كورسيهها} \textsuperscript{(32)}.

%%K t03-c1-14-1 p
14. --- شو ما أسعدها ! نسيت زهورها المفضّلة، \VUgras{القرنفل}
\textsuperscript{(34)} اللي عم يحطّ عليه \VUgras{الفراش}
\textsuperscript{(35)}، و\VUgras{التوليب} \textsuperscript{(36)}
و\VUgras{زنبق الوادي} \textsuperscript{(37)}، بـ\VUgras{قصاري}
\textsuperscript{(38)} عالـ\VUgras{حيط} \textsuperscript{(33)}،
و\VUgras{وردها} \textsuperscript{(39)} اللي عامل سياج كتير حلو. عم تتفرّج
على لبس السيّدات الغني بالموكب.

%%K t03-c1-15-1 p
15. --- ولاد الضيعة ما بينتبهوا كتير عاللبس. عم يتدافشوا ويزقّوا بعض
ليصيبوا \VUgras{ملبّس} \textsuperscript{(55)} أو \VUgras{مصاري}
\textsuperscript{(52)}. وواحد منهن عم يبكي \textsuperscript{(40)}.

%%K t03-c1-16-1 p
16. --- والتاني داس على إجر \VUgras{الكلب} \textsuperscript{(42)}، اللي
هرب وهوّي عم يعوي وهرّب \VUgras{الدجاجة} \textsuperscript{(43)}
و\VUgras{صيصانها} \textsuperscript{(44)}.

%%K t03-tit-6 sub
\VUsaut{5.0mm}
\VUpk{10.2pt}{40}{موكب العمادة.}
\VUsaut{4.0mm}

%%K t03-c1-17-1 p
17. --- قدّام الموكب عم تمشي \VUgras{المرضعة} \textsuperscript{(45)}.
عليها \VUgras{طاقية} \textsuperscript{(46)} حلوة مع شرايط طويلة. وشايلة
\VUgras{المولود} \textsuperscript{(47)} الملبّس كله \VUgras{دانتيل}
\textsuperscript{(48)}.

%%K t03-c1-18-1 p
18. --- ووراء المرضعة بيجي \VUgras{العرّاب} \textsuperscript{(49)}
و\VUgras{العرّابة} \textsuperscript{(53)}. هنّي عم يوزّعوا الملبّس
والمصاري. العرّاب عليه \VUgras{كفوف} \textsuperscript{(51)}
و\VUgras{برنيطة عالية} \textsuperscript{(50)}. والعرّابة ماسكة بإيدها
\VUgras{باقة ورد} \textsuperscript{(54)} حلوة.

%%K t03-c1-19-1 p
19. --- ووراهن، منشوف \VUgras{عمّ} \textsuperscript{(56)}
و\VUgras{عمّة} \textsuperscript{(58)} الولد. العمّة عليها
\VUgras{برنيطة} \textsuperscript{(59)} أنيقة، والعمّ عليه
\VUgras{ردنغوت} \textsuperscript{(57)} أسود و\VUgras{صدرية} بيضا. وهيّن
بعدين \VUgras{ابن العمّ} \textsuperscript{(60)} و\VUgras{بنت العمّ}
\textsuperscript{(61)}. وهاي ماسكة بإيدها \VUgras{أخت}
\textsuperscript{(62)} المولود، برتا الزغيرة.

%%K t03-c1-20-1 p
20. --- و\VUgras{أخو} \textsuperscript{(63)} برتا التاني عم يمشي وراء.
عم يعطي إيده لـ\VUgras{قريبة} \textsuperscript{(64)}. و\VUgras{رفقات}
\textsuperscript{(65)} العيلة بيجوا بعدين.

%%K t03-tit-7 sub
\VUsaut{4.6mm}
\VUpk{10.2pt}{40}{المتفرّجين.}
\VUsaut{3.4mm}

%%K t03-c1-21-1 p
21. --- جنب الموكب، هيّو \VUgras{شحّاد} \textsuperscript{(66)} متّكي على
\VUgras{عكّازه} \textsuperscript{(67)}. عم يطلب صدقة.

%%K t03-c1-22-1 p
22. --- وعالجهة التانية في صبي زغير كتير \VUgras{مؤدّب}
\textsuperscript{(68)}. عم يسلّم وبيقول « \VUgras{صباح الخير} » للناس
اللي بيعرفهن.

%%K t03-c1-23-1 p
23. --- أهل الضيعة طلعوا من بيوتهن ليتفرّجوا على موكب العمادة. منشوف
\VUgras{فلّاح} \textsuperscript{(69)} على راسه \VUgras{طاقية قطن}
وجنبه \VUgras{فلّاحة} \textsuperscript{(70)}. وبعدين \VUgras{ختيارة}
\textsuperscript{(71)} متّكية على \VUgras{عصاتها} \textsuperscript{(72)}.
وأخيراً \VUgras{الطحّان} \textsuperscript{(73)} مع \VUgras{برنيطة
اللبّاد} \textsuperscript{(74)} الكبيرة، وفلّاحة شابّة لابسة
\VUgras{قباقيب} \textsuperscript{(75)}، عم تعلّم ابنها الزغير يمشي.

%%K t03-c1-24-1 p
24. --- هالمشهد كتير مسلّي.

\VUsaut{4.0mm}
%%K t03-tit-8 sub
\VUcentre{12.0pt}{0}{\textit{المشهد التاني.}}
\VUsaut{3.0mm}

%%K t03-tit-9 sub
\VUpk{10.2pt}{40}{عيد عام.}
\VUsaut{2.4mm}

%%K t03-tit-10 sub
\VUpk{10.2pt}{40}{ألعاب مايّة وسباقات.}
\VUsaut{3.0mm}

%%K t03-c2-01-1 p
1. --- إجا \VUgras{الصيف}. شمس شهر \VUgras{حزيران} (تموز أو
\VUgras{آب}) الحامية بتشجّع الشباب عالسباحة والتجديف.

%%K t03-c2-02-1 p
2. --- عالضفّة اليمنى من النهر، المجدّفين عم يشلحوا تيابهن ليشاركوا
بالألعاب المايّة والسباقات.

%%K t03-c2-03-1 p
3. --- واحد منهن عم يشلح \VUgras{صبابيطه} \textsuperscript{(1)} ؛ عم
يفكّ رباطها. ورفيقيه التنين جاهزين من قبل. من هلّق ما عليهن غير
\VUgras{الكنزة} \textsuperscript{(2)} و\VUgras{لباس البحر}
\textsuperscript{(3)}. عم يحكوا وهنّي عم يستنّوا رفقاتهن. عم يحكوا عن
العيد والحفلة اللي عم تصير عالضفّة التانية.

%%K t03-c2-04-1 p
4. --- وعالأرض، جنبهن، مرميّة \VUgras{تياب} رفقاتهن : فردة
\VUgras{جزمات} \textsuperscript{(4)}، و\VUgras{صبابيط}
\textsuperscript{(5)}، و\VUgras{كلسات} \textsuperscript{(6)}،
وبرانيط \VUgras{لبّاد} \textsuperscript{(7)} و\VUgras{قش}
\textsuperscript{(8)}، و\VUgras{غرافة} \textsuperscript{(9)}.

%%K t03-c2-05-1 p
5. --- وواحد من الشباب طوى \VUgras{بدلته} \textsuperscript{(10)}
بعناية. حطّها على منشفة، وجاره عن قريب رح يلفّ فيها
\VUgras{البدلتين}.

%%K t03-c2-06-1 p
6. --- والسادس متأخّر. لسّا \VUgras{الكاسكيت} \textsuperscript{(11)}
على راسه. ما شلح لا \VUgras{قميصه} \textsuperscript{(12)} ولا
\VUgras{ياقته} \textsuperscript{(13)} ولا \VUgras{أساوره}
\textsuperscript{(14)}. وماسك بإيده \VUgras{صدريته}
\textsuperscript{(17)}، ولسّا عليه \VUgras{بنطلونه}
\textsuperscript{(16)} و\VUgras{حمّالاته} \textsuperscript{(15)}.
أكيد ما رح يكون جاهز للسباق الجاي.

%%K t03-c2-07-1 p
7. --- وجنب الشباب، في درب زغير عجوانبه \VUgras{شوك}
\textsuperscript{(18)} كتير، بيوصل لضفّة النهر، وين أبو يعقوب عم يحضّر
بين \VUgras{القصب} \textsuperscript{(19)} \VUgras{الألعاب النارية}
\textsuperscript{(20)} اللي رح يطيّروها بالليل.

%%K t03-tit-11 sub
\VUsaut{4.4mm}
\VUpk{10.2pt}{40}{السباق.}
\VUsaut{3.4mm}

%%K t03-c2-08-1 p
8. --- وعالنهر، كم سيّدة، مستظلّات بشمسياتهن، عم يتمشّوا
بـ\VUgras{قارب} \textsuperscript{(21)}. عم يتابعوا السباق اللي كتير
مشوّق. بكل \VUgras{زورق} \textsuperscript{(22)}، أربع \VUgras{مجدّفين}
\textsuperscript{(24)} عم يجدّفوا بكل قوّتهن، بينما \VUgras{الرُبّان}
\textsuperscript{(26)} ماسك \VUgras{الدفّة} \textsuperscript{(25)}
وبيوجّه القارب الرفيع. و\VUgras{المجاديف} \textsuperscript{(23)} عم
تضرب الميّ بإيقاع، والقوارب الخفيفة عم تسير بسرعة بتدوّخ.

%%K t03-c2-09-1 p
9. --- وكم شاب عم يتنافسوا، جنب عمود الجسر، على جايزة \VUgras{الصاري
المايل} \textsuperscript{(28)}. واحد منهن عم يمشي بحذر عالصاري المرن
والزحّاط ليوصل لـ\VUgras{العلم} \textsuperscript{(29)} الزغير المثبّت
بالطرف. و\VUgras{منافسه} \textsuperscript{(30)} قليل الحظ وقع
بالميّ. عم يسبح ليطلع عالبرّ.

%%K t03-c2-10-1 p
10. --- وشباب أكبر \VUgras{عم يتبارزوا عالميّ} \textsuperscript{(32)}
جنب \VUgras{الرصيف} \textsuperscript{(33)}. وكل هاللعب عم يتفرّج عليه
\VUgras{جمهور} \textsuperscript{(31)} كبير متّكي على \VUgras{درابزين}
\VUgras{الجسر} \textsuperscript{(27)} ومتجمّع عالضفّة. بس كم متفرّج عم
يلتفتوا ليتابعوا لعبة \VUgras{صاري الجايزة} \textsuperscript{(45)} أو
ليتفرّجوا على \VUgras{موكب البلدية} اللي جاي لـ\VUgras{توزيع}
الجوايز.

%%K t03-c2-11-1 p
11. --- أولاً، هيّن كم \VUgras{شقي} \textsuperscript{(35)} عم يمشوا
قدّام \VUgras{الموسيقيين} \textsuperscript{(36)} ؛ وبعدين بيجي
\VUgras{رئيس البلدية} \textsuperscript{(37)}، والمساعدين
و\VUgras{المجلس البلدي} تبع البلدة. وبالآخر عم تمشي
\VUgras{الإطفائية} \textsuperscript{(38)}.

%%K t03-tit-12 sub
\VUsaut{5.0mm}
\VUpk{10.2pt}{40}{العيد.}
\VUsaut{3.6mm}

%%K t03-c2-12-1 p
12. --- في ناس كتير على \VUgras{ساحة العيد} \textsuperscript{(39)}.
الناس عم تيجي تتفرّج عالـ\VUgras{بسطات} المختلفة : \VUgras{الخيل
الخشبية} \textsuperscript{(40)}، و\VUgras{السيرك} \textsuperscript{(41)}
اللي بلّش عنده العرض ؛ و\VUgras{الحفلة الراقصة} \textsuperscript{(42)}،
وين شباب وصبايا البلدة عم يتمتّعوا بلذّة \VUgras{الرقص}.

%%K t03-c2-13-1 p
13. --- وبالآخر، منشوف \VUgras{الحلبة} \textsuperscript{(44)} وين
\VUgras{المصارع} الإسباني الشجاع عم يصارع \VUgras{التور}
\textsuperscript{(45)} الهايج، وبعدين \VUgras{الزحليقة}
\textsuperscript{(49)}، و\VUgras{ميدان الرماية} \textsuperscript{(46)}
وين الناس عم تتمرّن تستعمل \VUgras{البواريد} وتضرب عالـ\VUgras{هدف}.

%%K t03-c2-14-1 p
14. --- وعن قريب، هيّو \VUgras{مسرح العيد} \textsuperscript{(47)} اللي
بيمثّلوا فيه \VUgras{الكوميديا} و\VUgras{الدراما}. وكتير قريب، في
بسطة \VUgras{العملاق} \textsuperscript{(48)} و\VUgras{القزم}.
\VUgras{طول} الأول متران وخمستعشر سنتيمتر ؛ والتاني بالكاد بطول ولد.

%%K t03-c2-15-1 p
15. --- وأخيراً، هيّي \VUgras{حديقة الحيوانات} \textsuperscript{(50)}
الكبيرة مع \VUgras{وحوشها} : \VUgras{النمر} \textsuperscript{(51)}
بـ\VUgras{مخالبه} القوية، و\VUgras{الأسد} \textsuperscript{(52)}
بـ\VUgras{لبدته} الكتيفة، و\VUgras{زرافتين} \textsuperscript{(53)}
و\VUgras{الفيل} \textsuperscript{(54)} اللي بيعرف يستعمل
\VUgras{خرطومه} بمهارة.

%%K t03-c2-16-1 p
16. --- و\VUgras{المروّض} \textsuperscript{(55)} اللي عم يدرّب هالوحوش
واقف عند باب البسطة.

\VUsaut{6.0mm}
\VUfilet{22mm}
